%! Scatterplots and Correlation — Unit 2, Lesson 2.2 — Warm-Up (Answer Key)
% ONE PAGE, blank and key. Three short items.
% Item 1 spirals to Lesson 2.1 (association between two categorical variables, by comparing
% conditional percents) and plants the word ASSOCIATION. Item 2 spirals to Unit 1 (the four
% things every description of a distribution names) and seeds DOFS, today's four-part
% description. Item 3 seeds explanatory/response and direction with five pairs of numbers.
% Mirrors the blank exactly apart from the package swap, the pageheader, and the answers.
% NO teachernote here — teacher prose lives in the lesson plan.
\documentclass[12pt]{article}
\usepackage{apstats-article}
\usepackage{apstats-key}   % pulls in -boxes; do NOT also load -boxes

\begin{document}

\pageheader{Unit 2, Lesson 2.2}{Warm-Up --- Answer Key}
% NAMESTRIP: no \namedateperiod here — the name row lives on the cover only.

\vspace{0.28in}

\begin{enumerate}[leftmargin=1.9em, itemsep=0.75in]

  \item \textbf{From last lesson.} At one school, 60 of the 80 students who play a sport
        passed a fitness test, and 72 of the 120 who do not play a sport passed.

        \vspace{12pt}
        Pass rate, sport: \ans{75\%} \qquad no sport: \ans{60\%}
        \vspace{16pt}

        Is there an \emph{association} between playing a sport and passing? How can you
        tell?\\[18pt]
        \ansline{Yes --- the pass rates differ (75\% vs.\ 60\%),}\\[1pt]
        \ansline{so knowing one variable helps predict the other.}

  \item \textbf{From Unit 1.} Every description of a distribution of one quantitative
        variable names four things. List them.

        \vspace{16pt}
        \ans{shape} \quad \ans{center} \quad \ans{spread} \quad
        \ans{outliers/unusual}

  \item Five students reported hours studied and their quiz score.

        \vspace{12pt}
        \renewcommand{\arraystretch}{1.5}
        \begin{tabularx}{\linewidth}{@{} l Y Y Y Y Y @{}}
        \rowcolor{sky}
        \textbf{Hours studied} & 1 & 2 & 3 & 4 & 5 \\
        \hline
        \textbf{Quiz score} & 62 & 70 & 74 & 81 & 88 \\
        \end{tabularx}

        \vspace{16pt}
        As hours studied go up, what tends to happen to the score? \ans{it tends to go up}
        \vspace{16pt}

        Which variable would you use to \emph{predict} the other? Why?\\[18pt]
        \ansline{Hours studied: it comes first, and it could help}\\[1pt]
        \ansline{explain or predict the score.}

\end{enumerate}

\end{document}
